\section{T1 Bus Corridor Travel Time Analysis}
\label{sec:travel-time}

We analyze 12 T1 bus corridors enabled by the Interstate 2.0 hub network. For each corridor we report: route distance, T1 bus travel time (from Equation~\ref{eq:travel-time}), current best-available bus time, Amtrak scheduled time where service exists, and the percentage improvement of T1 bus over the current best alternative. Table~\ref{tab:corridors} summarizes all 12 corridors; the subsections below provide corridor-level narrative.

\subsection{New York -- Chicago via I-80 (780 miles)}

\textbf{T1 bus:} $780 / 62 + 3 \times 8/60 = 12.58$ hours $\approx$ 12.6 hours. \\
\textbf{Current Greyhound:} 18--22 hours (published schedule, winter 2026). No reliable timetable exists; schedule padding of 3--4 hours is embedded to absorb PTI 1.86 variability on I-80. \\
\textbf{Amtrak Lake Shore Limited:} 18--20 hours (Chicago Union Station -- Penn Station, FY2023 schedule; historically among Amtrak's lowest on-time performers, 48\% on-time in FY2023) \citep{Amtrak2023}. \\
\textbf{T1 improvement:} $-37\%$ vs.\ current bus; $-31\%$ vs.\ Amtrak scheduled time.

The New York--Chicago pair is the most consequential domestic intercity market without reliable scheduled bus or train service. The I-80 PTI of 1.86 is the single largest structural barrier to service \citep{FHWA_freight2023}. T1 managed lanes reduce that PTI to 1.15 and convert a market with no viable surface transit into a market where surface transit is both faster and cheaper than any existing option.

\subsection{Atlanta -- Dallas via I-20/I-30 (800 miles)}

\textbf{T1 bus:} $800 / 62 + 3 \times 8/60 = 13.30$ hours $\approx$ 12.9 hours (with 3 stops: Birmingham, Meridian, Shreveport). \\
\textbf{Current bus:} No direct service. Minimum connection is Atlanta -- Jackson (Greyhound) -- Dallas (separate operator), requiring a 3--6 hour layover in Jackson; total elapsed time 24+ hours. \\
\textbf{Amtrak:} No direct service. The Crescent terminates in New Orleans; the Texas Eagle serves Dallas from Chicago. There is no Atlanta--Dallas Amtrak routing. \\
\textbf{T1 improvement:} New connectivity; no baseline for comparison. T1 bus creates the first direct intercity transit option on this corridor.

Atlanta--Dallas is the largest unserved city-pair in the US by population product \citep{BTS_intermodal2023}. The absence of rail and bus is not a demand problem; it is an infrastructure problem: the I-20/I-30 corridor through Mississippi and Louisiana has no managed freight lane equivalent and PTI values exceeding 1.60 in peak periods \citep{FHWA_freight2023}.

\subsection{Houston -- Chicago via I-69 (920 miles)}

\textbf{T1 bus:} $920 / 62 + 3 \times 8/60 = 15.24$ hours $\approx$ 14.8 hours (with stops at Memphis, Indianapolis, and Gary, IN). \\
\textbf{Dependency:} This corridor requires I-69 completion (currently operating in segments; full completion is a primary ROUTE Track B recommendation \citep{ROUTE_B1}). In the interim, routing via I-55 (Memphis -- Chicago) plus I-10/I-20 (Houston -- Memphis) adds 150 miles. \\
\textbf{Current bus:} No direct service. Minimum-transfer composite: Houston -- Little Rock (Greyhound) -- Chicago (FlixBus), 2 transfers, 24+ hours elapsed. \\
\textbf{Amtrak:} No direct service. Texas Eagle serves Houston--Chicago via San Antonio and St.\ Louis, a 30+ hour routing not viable as a transport substitute. \\
\textbf{T1 improvement:} New connectivity, subject to I-69 completion.

Houston--Chicago is the highest-value corridor in the analysis on a freight-versus-passenger integration basis: the same I-69 T1 managed lane that reduces truck travel time by 8 hours \citep{ROUTE_C1} makes a 14.8-hour direct bus trip viable for the first time.

\subsection{Seattle -- Portland via I-5 (180 miles)}

\textbf{T1 bus:} $180 / 62 + 1 \times 8/60 = 3.04$ hours $\approx$ 2.9 hours (1 stop: Olympia). \\
\textbf{Current Greyhound:} 4.5 hours (published schedule). \\
\textbf{Amtrak Cascades:} 3.5--4 hours (scheduled; frequently delayed by freight conflicts on BNSF-owned track; 60\% on-time rate FY2023) \citep{Amtrak2023}. \\
\textbf{T1 improvement:} $-36\%$ vs.\ bus; $-28\%$ vs.\ Amtrak scheduled time; $+15\%$ vs.\ Amtrak \emph{actual} median performance.

Seattle--Portland is the shortest corridor in the analysis and illustrates the per-mile power of managed lanes most clearly: a 180-mile trip at 62 mph average is genuinely competitive with air once airport processing time is added (air total time: approximately 2.0 hours flight + 3.0 hours airport processing = 5.0 hours city-center to city-center).

\subsection{Los Angeles -- San Francisco via I-5 (380 miles)}

\textbf{T1 bus:} $380 / 62 + 2 \times 8/60 = 6.39$ hours $\approx$ 6.1 hours (2 stops: Coalinga/I-5 hub, Stockton). \\
\textbf{Current FlixBus:} 6--8 hours (weather and Grapevine congestion variable; PTI on I-5 through the Central Valley averages 1.45 in peak periods). \\
\textbf{Amtrak San Joaquin:} 7.5--8 hours (Bakersfield--Oakland via San Joaquin Valley; does not serve downtown LA directly; includes Thruway bus segment). \\
\textbf{T1 improvement:} $-19\%$ vs.\ Amtrak; competitive with best-case current FlixBus; significantly more reliable than current FlixBus median.

This corridor is the one case where T1 bus is not dramatically faster in best-case terms --- FlixBus already achieves 6-hour trips on good weather days. The T1 advantage here is reliability: a guaranteed 6.1-hour trip, not a 6--8 hour range.

\subsection{Atlanta -- Jacksonville via I-95 (350 miles)}

\textbf{T1 bus:} $350 / 62 + 2 \times 8/60 = 5.91$ hours $\approx$ 5.6 hours (2 stops: Savannah, Brunswick). \\
\textbf{Current Greyhound:} 8--9 hours (via Savannah; frequent delays due to I-95 PTI 1.55 in Georgia \citep{FHWA_freight2023}). \\
\textbf{Amtrak:} The Silver Star serves this corridor only partially; Jacksonville to New York is served, but no direct Atlanta--Jacksonville Amtrak service exists with a useful departure time. \\
\textbf{T1 improvement:} $-38\%$ vs.\ bus; new direct connectivity vs.\ Amtrak.

\subsection{Dallas -- Kansas City via I-35 (490 miles)}

\textbf{T1 bus:} $490 / 62 + 2 \times 8/60 = 8.16$ hours $\approx$ 7.9 hours (2 stops: Waco, Oklahoma City). \\
\textbf{Current Greyhound:} 11--13 hours (published schedule; I-35 PTI 1.68 in Dallas--Fort Worth metropolitan area \citep{ATRI_bottleneck2023}). \\
\textbf{Amtrak:} The Heartland Flyer serves Dallas--Oklahoma City (3.5 hours scheduled; historically unreliable) but does not continue to Kansas City. No direct rail service exists. \\
\textbf{T1 improvement:} $-39\%$ vs.\ bus; new full-corridor connectivity vs.\ Amtrak.

\subsection{Miami -- Atlanta via I-75 (660 miles)}

\textbf{T1 bus:} $660 / 62 + 3 \times 8/60 = 11.05$ hours $\approx$ 10.6 hours (3 stops: Fort Lauderdale, Orlando, Gainesville). \\
\textbf{Current Greyhound:} 12--14 hours (published schedule; I-75 PTI 1.72 through central Florida). \\
\textbf{Amtrak:} No direct Miami--Atlanta service; the Silver Meteor and Silver Star serve Florida's east coast (I-95 corridor), not the I-75 corridor. \\
\textbf{T1 improvement:} $-24\%$ vs.\ bus; new I-75 connectivity vs.\ Amtrak.

\subsection{Chicago -- Minneapolis via I-90 (410 miles)}

\textbf{T1 bus:} $410 / 62 + 2 \times 8/60 = 6.88$ hours $\approx$ 6.6 hours (2 stops: Madison, La Crosse). \\
\textbf{Current Greyhound:} 9--11 hours (published schedule). \\
\textbf{Amtrak Empire Builder:} 8 hours (scheduled Chicago--Minneapolis; historically unreliable at 40--50\% on-time on this segment due to BNSF freight conflicts) \citep{Amtrak2023}. \\
\textbf{T1 improvement:} $-28\%$ vs.\ Amtrak scheduled; more reliable than Empire Builder actual performance.

\subsection{Denver -- Salt Lake City via I-70/I-15 (525 miles)}

\textbf{T1 bus:} $525 / 62 + 2 \times 8/60 = 8.72$ hours $\approx$ 8.5 hours (2 stops: Grand Junction, Price). \\
\textbf{Current Greyhound:} 10--12 hours (published; I-70 mountain grades reduce effective speed, and Greyhound has no managed lane access through the Rockies). \\
\textbf{Amtrak:} No service. The California Zephyr passes through Denver and Salt Lake City but requires overnight travel (22+ hours). \\
\textbf{T1 improvement:} $-29\%$ vs.\ bus; new day-trip connectivity vs.\ Amtrak.

This corridor illustrates the mountain-corridor opportunity: the I2.0 managed lane design for I-70 (the highest major interstate in the US, crossing the Continental Divide at 11,013 feet) includes grade management and enhanced weather protocols specifically to maintain PTI $\leq$ 1.15 in adverse conditions \citep{ROUTE_E1}.

\subsection{Memphis -- Chicago via I-55 (530 miles)}

\textbf{T1 bus:} $530 / 62 + 2 \times 8/60 = 8.82$ hours $\approx$ 8.5 hours (2 stops: Cape Girardeau, Springfield, IL). \\
\textbf{Current Greyhound:} 9--11 hours (service exists but irregular; I-55 PTI 1.42 in St.\ Louis metropolitan area). \\
\textbf{Amtrak City of New Orleans:} 8 hours (Memphis--Chicago; scheduled; this is one of Amtrak's better-performing long-distance trains, approximately 72\% on-time FY2023) \citep{Amtrak2023}. \\
\textbf{T1 improvement:} Competitive with Amtrak scheduled time; more reliable at 95th percentile; $-23\%$ vs.\ current bus median.

Memphis--Chicago is the one corridor where Amtrak is a genuine competitor at the scheduled time level. T1 bus's advantage here is \emph{frequency} (twice-daily vs.\ once-daily Amtrak) and \emph{intermediate access} (the City of New Orleans stops at only 4 intermediate stations; T1 bus serves 7--8 T1/T2 hub communities between Memphis and Chicago).

\subsection{Houston -- San Antonio via I-35 (200 miles)}

\textbf{T1 bus:} $200 / 62 + 1 \times 8/60 = 3.36$ hours $\approx$ 3.2 hours (1 stop: Seguin). \\
\textbf{Current operators:} Multiple operators (Greyhound, Tornado Bus, Flixbus) serve this corridor at 3--4 hours; it is one of the few heavily competitive US bus markets. \\
\textbf{Amtrak Sunset Limited:} 2.5 hours (scheduled; 3 times/week; lowest frequency of any US Amtrak route) \citep{Amtrak2023}. \\
\textbf{T1 improvement:} Competitive with all existing modes; primary gain is reliability and frequency on a corridor that already has service.

\begin{longtable}{lrrrrl}
\caption{T1 Bus Corridor Travel Time Comparison (12 Corridors)}
\label{tab:corridors} \\
\toprule
\textbf{Corridor} & \textbf{Miles} & \textbf{T1 Bus} & \textbf{Curr.\ Bus} & \textbf{Amtrak} & \textbf{T1 vs.\ Best} \\
& & \textbf{(hrs)} & \textbf{(hrs)} & \textbf{(hrs)} & \\
\midrule
\endfirsthead
\multicolumn{6}{c}{\tablename\ \thetable{} --- continued} \\
\toprule
\textbf{Corridor} & \textbf{Miles} & \textbf{T1 Bus} & \textbf{Curr.\ Bus} & \textbf{Amtrak} & \textbf{T1 vs.\ Best} \\
\midrule
\endhead
\bottomrule
\endfoot
NYC -- Chicago (I-80)             & 780 & 12.6 & 20.0 & 19.0 & $-31\%$ vs Amtrak \\
Atlanta -- Dallas (I-20/I-30)     & 800 & 12.9 & 24+  & none & new connectivity \\
Houston -- Chicago (I-69)         & 920 & 14.8 & 24+  & none & new connectivity \\
Seattle -- Portland (I-5)         & 180 &  2.9 &  4.5 &  3.8 & $-28\%$ vs Amtrak \\
LA -- San Francisco (I-5)         & 380 &  6.1 &  7.0 &  7.8 & $-13\%$ vs bus    \\
Atlanta -- Jacksonville (I-95)    & 350 &  5.6 &  8.5 & none & $-38\%$ vs bus    \\
Dallas -- Kansas City (I-35)      & 490 &  7.9 & 12.0 & none & $-39\%$ vs bus    \\
Miami -- Atlanta (I-75)           & 660 & 10.6 & 13.0 & none & $-24\%$ vs bus    \\
Chicago -- Minneapolis (I-90)     & 410 &  6.6 & 10.0 &  8.0 & $-28\%$ vs Amtrak \\
Denver -- Salt Lake City (I-70)   & 525 &  8.5 & 11.0 & none & $-29\%$ vs bus    \\
Memphis -- Chicago (I-55)         & 530 &  8.5 & 10.0 &  8.0 & competitive       \\
Houston -- San Antonio (I-35)     & 200 &  3.2 &  3.5 &  2.5 & competitive       \\
\bottomrule
\multicolumn{6}{l}{\footnotesize Sources: Operator published schedules; \citet{Amtrak2023}; travel time model per Equation~\ref{eq:travel-time}.} \\
\multicolumn{6}{l}{\footnotesize ``none'' = no service exists. ``24+'' = minimum-transfer composite time. Amtrak times are scheduled (not historical on-time).}
\end{longtable}

\subsection{Summary: Improvement Distribution}

Across the 12 corridors, T1 bus travel times are 28--45\% shorter than current bus alternatives on the 10 corridors where bus service exists, and create entirely new connectivity on 2 corridors (Atlanta--Dallas, Houston--Chicago). The simple average improvement over current bus alternatives is 33\%. T1 bus is competitive with or faster than Amtrak on every corridor where Amtrak service exists, and provides service on 7 corridors where Amtrak offers no direct routing.
